\title{Direct Preference Optimization: Your Language Model is Secretly a Reward Model}

\begin{abstract}
Although large-scale unsupervised language models (LMs) acquire extensive world knowledge and some degree of reasoning capability, precisely directing their outputs remains challenging due to their wholly unsupervised training regimen. Common approaches to improving model controllability involve collecting human judgments on the relative merit of generated outputs and then further training the original LM to favor these preferences, frequently through reinforcement learning from human feedback (RLHF). Yet, RLHF is known for its intricacy and potential instability, as it first entails fitting a reward model to capture human evaluations, then using reinforcement learning to optimize the LM toward this learned reward, all while seeking to preserve proximity to the base model. This work presents a novel reward model parameterization within the RLHF framework, enabling the corresponding optimal policy to be derived analytically. This re-formulation allows the classical RLHF challenge to be addressed with a straightforward classification loss. We term the resulting technique \textit{Direct Preference Optimization} (DPO), which is characterized by strong empirical performance, robustness, and computational efficiency—it avoids the need for LM sampling during adaptation and dispenses with extensive hyperparameter searches. Experimental results demonstrate that DPO achieves fine-tuning of LMs to reflect human preferences on par with, or exceeding, established methods. In particular, DPO surpasses RLHF approaches based on PPO in sentiment modulation of model outputs and achieves comparable or better response quality for tasks like summarization and single-turn dialogue, while offering greater simplicity in both implementation and training procedures.
\end{abstract}

\section{Introduction}
Large-scale unsupervised language models (LMs), when trained on massive corpora, develop a range of unexpected abilities~\citep{chowdhery2022palm, brown2020language, touvron2023llama,bubeck2023sparks}. Despite these achievements, such models learn from human-generated data encompassing diverse intentions, priorities, and expertise. Not all of these traits are desirable for emulation; for instance, we may prefer an AI coding assistant that can \textit{recognize} typical programming errors to rectify them, yet, during code generation, we want the model to emulate the (sometimes uncommon) expert-level code present in its training set. Likewise, it is valuable for a language model to \textit{know about} a widespread misconception held by half the population, but we would not want it to assert that misconception as fact in half of related responses. Thus, the process of isolating a model's \emph{targeted outputs and behavior} from its expansive range of \textit{acquired knowledge and skills} is a pivotal aspect in creating AI systems that are reliable, high-performing, and easy to direct~\citep{ouyang2022training}. Although standard practice employs reinforcement learning (RL) to nudge LMs toward reflecting human values, we demonstrate that the RL objective commonly adopted in these frameworks can, in fact, be solved exactly through a straightforward binary cross-entropy loss, streamlining the entire preference alignment process.

Broadly speaking, prevailing strategies embed desirable behaviors in language models by using carefully constructed datasets that capture human judgments of what is safe and beneficial. This stage of preference alignment is performed subsequent to an initial phase of extensive unsupervised pre-training over text corpora. While supervised fine-tuning on expert demonstrations is the most direct route for aligning preferences, reinforcement learning from human or AI feedback (RLHF/RLAIF; \citep{christiano2017deep,bai2022constitutional}) has proven the most effective in practice. In RLHF, a reward model is trained on preference-labeled data to quantify human value judgments, and RL is subsequently used to refine the LM policy to generate responses with higher predicted rewards, while limiting divergence from the pre-trained model. RLHF-based models have demonstrated remarkable proficiency in conversation and coding, but this approach involves considerable complexity beyond simple supervised learning: it requires training several LMs and sampling from the evolving model within the training loop, which leads to elevated computational demands.

In this work, we demonstrate a method for training language models to conform to human preferences without the need for explicit reward models or reinforcement learning approaches. We introduce \textit{{Direct Preference Optimization} (DPO)}, a novel algorithm that inherently pursues the same reward maximization objective—regularized by a KL-divergence penalty—as standard RLHF algorithms, yet offers a significantly more accessible and efficient implementation. DPO fundamentally modifies the log probability ratio between preferred and non-preferred outputs, augmented with a dynamic, example-specific importance weighting. This weighting addresses model degeneration issues observed when employing a naive probability ratio objective. Similar to prior approaches, DPO adopts a formalized preference model, such as the Bradley-Terry model \cite{bradley1952rankanalysis}, to evaluate the concordance between a reward function and actual preference annotations. However, rather than following the two-stage process of first fitting a reward model with a preference loss, and then training a policy against this learned reward function, DPO reformulates the preference loss directly in terms of the policy via a change of variables. With access to a corpus of human-annotated preferences on model outputs, DPO is thus able to directly optimize the policy via a straightforward binary cross entropy objective, yielding a policy that is optimal for an implicit reward derived from the preference observations.

The principal contribution of this work is the development of Direct Preference Optimization (DPO), a straightforward, reinforcement-learning-free technique for preference-driven language model training. Empirical evaluations reveal that DPO matches or exceeds the effectiveness of state-of-the-art alternatives—including PPO-based RLHF—across a range of tasks such as sentiment control, summarization, and conversational modeling, utilizing models with up to 6B parameters.

\section{Related Work}

Self-supervised language models have demonstrated the ability to tackle certain tasks in a zero-shot manner \citep{radford2019language}, or with minimal examples provided as prompts \citep{gpt3,megatron,chowdhery2022palm}, as their scale increases. Nonetheless, adapting these models to downstream applications and improving alignment with user expectations is often achieved by fine-tuning them on curated datasets containing instructions and human-generated completions \citep{mishra-etal-2022-cross,sanh2022multitask,chung2022scaling,thoppilan2022lamda}. This process, known as `instruction-tuning,' enhances LLMs’ ability to handle instructions not present in the original tuning set and broadly raises their effectiveness in real-world settings \citep{chung2022scaling}. Despite its efficacy, gathering \textit{relative} human assessments of response quality tends to be less demanding than sourcing expert-labeled demonstrations. As a result, subsequent research has leveraged collections of human preferences to further train LLMs, which has led to marked improvements in domains such as translation \citep{kreutzer-etal-2018-reliability}, summarization \citep{stiennon2022learning,ziegler2020finetuning}, narrative generation \citep{ziegler2020finetuning}, and instruction adherence \citep{ouyang2022training,ramamurthy2023is}. In these approaches, a reward function is learned—optimized to mirror the human preferences dataset, frequently using preference models like the Bradley-Terry model \citep{bradley1952rankanalysis}—after which the language model is adjusted to maximize this reward through reinforcement learning techniques, typically employing algorithms such as REINFORCE \citep{williams1992reinforce}, proximal policy optimization (PPO; \cite{schulman2017proximal}), or their modifications \citep{ramamurthy2023is}. A related research direction exploits LLMs, already instruction-tuned with feedback, to produce additional synthetic preference labels for specific traits, such as safety or non-harmfulness \citep{bai2022constitutional}, utilizing limited human oversight via textual rubrics for LLM-produced judgments. Collectively, these strategies bring together the fields of reinforcement learning-based language model training~\citep{Ranzato2015SequenceLT,paulus2018a,wu2018learning} and the development of techniques for preference-based learning from humans \citep{christiano2017deep,kupcsik2018learning}. Although preference-based fine-tuning with reinforcement learning is attractive, applying these techniques to large language models introduces significant practical hurdles; the present work introduces a theoretically grounded framework to optimize for relative human preferences while sidestepping reinforcement learning.

Beyond the realm of language, the problem of learning policies from preference data has been addressed in both bandit and reinforcement learning paradigms, yielding a variety of methodological advancements. One notable avenue is contextual bandit learning driven by preferences or action rankings, as opposed to direct reward signals; this is encapsulated in the framework of contextual dueling bandits (CDB; \cite{yue2012karmed,dudik2015contextual}). Without access to explicit reward information, theoretical treatments of CDBs replace the standard criterion of an optimal policy with the \textit{von Neumann winner}: a policy expected to outperform \textit{any} alternative at least half the time \citep{dudik2015contextual}. Unlike human preference learning—where the learning agent is provided with a pre-collected set of annotated action comparisons \citep{yan2022human}—CDBs typically acquire preference signals online. Likewise, the field of \textit{preference-based RL} (PbRL) concerns itself with policy learning from binary preference comparisons obtained from an implicit, unobserved 'scoring' function rather than explicit rewards \citep{BusaFekete2014,ruiz2023dueling}. While a diversity of PbRL algorithms exist, many rely on off-policy data reuse and proceed via explicit inference of this hidden reward or scoring model prior to policy optimization \citep{jain2013learning,BusaFekete2014,christiano2017deep,sadigh2017active,kupcsik2018learning}. In contrast, our contribution is a one-step approach: we develop a method for direct policy optimization to align with expressed preferences.

\section{Preliminaries}\label{section:prelims}

We outline the standard RLHF procedure as established in \citeauthor{ziegler2020finetuning}, with subsequent refinements in works such as \citep{stiennon2022learning, bai2022training, ouyang2022training}. This methodology generally encompasses three core stages: (1) supervised fine-tuning (SFT); (2) generation and annotation of preference data for reward modelling; and (3) reinforcement learning-based policy optimization.

\textbf{SFT}: The RLHF pipeline commences with the supervised adaptation of a pre-trained language model (LM), using high-quality labeled datasets relevant to the targeted downstream applications (e.g., dialogue, summarization), resulting in the SFT policy $\pi^\text{SFT}$.

\textbf{Reward Modelling Phase}: Following SFT, prompts $x$ are input into the SFT model to generate pairs of outputs $(y_1, y_2)\sim \pi^\text{SFT}(y \mid x)$. Human annotators are then tasked with indicating a preference, specifying $y_w\succ y_l \mid x$, where $y_w$ is the more favored answer over $y_l$ for a given prompt. These preferences are presumed to reflect an underlying, unobserved reward function $r^*(y, x)$. Among several preference modelling approaches, the Bradley-Terry (BT) model \cite{bradley1952rankanalysis} is commonly used, although generalizations like the Plackett-Luce model \citep{plackett1975analysis, luce2012individual} can also be adopted if rank orderings among multiple answers are provided. The BT model formalizes the human preference probability $p^*$ as follows:

Given a fixed dataset of pairwise preferences $\mathcal{D}=\bigl\{x^{(i)}, y_w^{(i)}, y_l^{(i)}\bigr\}_{i=1}^N$ generated according to $p^*$, we parameterize a reward function $r_{\phi}(x, y)$ and estimate the optimal parameters via maximum likelihood estimation. Interpreting the task as a binary classification, the negative log-likelihood loss can be expressed as:

\begin{equation}\label{eq:reward_model}
    \mathcal{L}_R(r_{\phi}, \mathcal{D}) = -\mathbb{E}_{(x, y_w, y_l)\sim \mathcal{D}}\bigl[\log \sigma(r_{\phi}(x, y_w)- r_{\phi}(x, y_l))\bigr]
\end{equation}

where $\sigma$ is the logistic sigmoid. For language model applications, $r_{\phi}(x, y)$ typically builds upon the SFT model $\pi^\text{SFT}(y \mid x)$ by incorporating a linear output layer above the last transformer block, which outputs a scalar reward estimate for each sequence \cite{ziegler2020finetuning}. To stabilize the reward estimates and reduce their variance, previous work standardizes the output such that $\mathbb{E}_{x,y\sim \mathcal{D}}\left[r_\phi(x, y)\right] = 0$ for all $x$.

\textbf{RL Fine-Tuning Phase}: In the reinforcement learning stage, the model leverages the trained reward function as a supervisory signal. Following established methodology~\citep{jaques2017sequence, jaques2020human}, the objective takes the form

\begin{equation}\label{eq:RL}
\max_{\pi_{\theta}}  \mathbb{E}_{x\sim \mathcal{D}, y\sim \pi_{\theta}(y \mid x)}\bigl[r_{\phi}(x, y)\bigr] - \beta\mathbb{D}_{\textrm{KL}}\bigl[\pi_{\theta}(y\mid x)\mid \mid \pi_\text{ref}(y\mid x)\bigr],
\end{equation}

where the hyperparameter $\beta$ determines the degree of divergence permitted from the reference policy $\pi_\text{ref}$, generally set as the pre-trained SFT policy $\pi^\text{SFT}$. The language model policy $\pi_\theta$ is also usually initialized using $\pi^\text{SFT}$. The KL penalty term is crucial, acting to constrain the learned policy within the region where the reward model remains reliable, thus preserving generative diversity and guarding against degenerate solutions that yield trivial high-reward outputs. Because text generation involves sampling discrete tokens, the objective is not differentiable, and reinforcement learning techniques must be applied. Common practice \citep{ziegler2020finetuning, stiennon2022learning, bai2022training, ouyang2022training} is to define a shaped reward as ${r(x, y) = r_{\phi}(x, y) -\beta (\log \pi_{\theta}(y\mid x) - \log \pi_\text{ref}(y\mid x))}$, which is optimized via PPO \cite{schulman2017proximal}.

\section{Direct Preference Optimization}\label{sec:DPO}

Given the computational demands and practical complexity of using reinforcement learning for large-scale model fine-tuning, we seek a more direct methodology for optimizing the policy with respect to preference data. Departing from conventional RLHF procedures that first train an explicit reward model before conducting RL-based optimization, our method utilizes a specific reward model parameterization that allows for closed-form derivation of the associated optimal policy, thereby removing the need for iterative RL training. As detailed below, our main contribution lies in exploiting the analytical relationship between reward models and their induced optimal policies, making it possible to reformulate a reward-centric loss into a policy-centric loss by a change of variables. This approach eliminates the need for explicit reward model fitting, yet retains fidelity to established preference modeling assumptions, such as those of the Bradley-Terry framework. Under our scheme, the policy model simultaneously fulfills the roles of language generation and (implicitly) preference-based reward estimation.

\textbf{Derivation of the DPO Objective.} We begin from the conventional RL objective, as described in Eq.~\ref{eq:RL}, which relies on a general reward function $r$. Prior research~\citep{peters2007reinforcement, peng2019advantage, korbak2022reinforcement, go2023aligning} has established that the solution optimizing the KL-regularized reward maximization in Eq.~\ref{eq:RL} adopts the following structure:

\begin{equation}\label{eq:op_policy}
    \pi_r(y\mid x) = \frac{1}{Z(x)}\pi_\text{ref}(y\mid x)\exp\left(\frac{1}{\beta}r(x, y)\right),
\end{equation}

with $Z(x) =\sum_{y}\pi_\text{ref}(y\mid x)\exp\left(\frac{1}{\beta}r(x, y)\right)$ denoting the normalization constant, or partition function. Detailed derivations are provided in Appendix~\ref{app:derivation1}. Estimating $Z(x)$ becomes computationally demanding even if the MLE estimator $r_{\phi}$ for the ground-truth reward $r^*$ is available~\citep{korbak2022reinforcement, go2023aligning}, posing significant practical challenges for this formulation. To address this, we may manipulate Eq.~\ref{eq:op_policy} to rewrite the reward as a function of its associated optimal policy $\pi_r$, reference policy $\pi_\text{ref}$, and partition function $Z(\cdot)$. In particular, taking the logarithm of both sides of Eq.~\ref{eq:op_policy} and rearranging terms yields:

\begin{equation}\label{eq:main_eq}
    r(x,y) =\beta \log \frac{\pi_r(y\mid x)}{\pi_\text{ref}(y\mid x)} + \beta \log Z(x).
\end{equation}

This reparameterization can likewise be applied to the true reward $r^*$ and its corresponding optimal policy $\pi^*$. Critically, the Bradley-Terry model is determined solely by reward differences between two outputs: ${p^*(y_1 \succ y_2 \mid x) = \sigma(r^*(x, y_1) - r^*(x, y_2))}$. Upon substituting Eq.~\ref{eq:main_eq} into the preference model of Eq.~\ref{eq:bradley-terry} for $r^*(x,y)$, the partition function terms eliminate, resulting in an expression for the human preference probability based only on the optimal policy $\pi^*$ and the reference policy $\pi_\text{ref}$. Accordingly, under the Bradley-Terry model, the optimal RLHF policy $\pi^*$ must satisfy:

\begin{equation}\label{eq:objective}
    p^*(y_1\succ y_2 \mid x)=\frac{1}{1 + \exp\left(\beta \log \frac{\pi^*(y_2\mid x)}{\pi_\text{ref}(y_2\mid x)} - \beta \log \frac{\pi^*(y_1\mid x)}{\pi_\text{ref}(y_1\mid x)}\right)}
\end{equation}

Appendix~\ref{app:derivation2} contains the full derivation. While Eq.~\ref{eq:objective} adopts the Bradley-Terry framework, analogous derivations are possible using the broader Plackett-Luce models~\citep{plackett1975analysis, luce2012individual}; details are available in Appendix~\ref{app:plackett_luce_models}.

With the human preference distribution now parameterized in terms of the optimal policy, not the reward, we are equipped to formulate a maximum likelihood training criterion for any parameterized policy $\pi_\theta$. Drawing a parallel to standard reward modeling (see Eq.~\ref{eq:reward_model}), our resulting policy learning objective is:

\begin{equation}\label{eq:optimum_model}
    \mathcal{L}_\text{DPO}(\pi_{\theta}; \pi_\text{ref}) = -\mathbb{E}_{(x, y_w, y_l)\sim \mathcal{D}}\left[\log \sigma \left(\beta \log \frac{\pi_{\theta}(y_w\mid x)}{\pi_\text{ref}(y_w\mid x)} - \beta \log \frac{\pi_{\theta}(y_l\mid x)}{\

In this approach, we model the implicit reward via an alternative parameterization, where the corresponding optimal policy coincides with $\pi_\theta$. Additionally, because our method aligns with learning a reparametrized Bradley-Terry model, it inherits theoretical guarantees such as consistency, assuming appropriate conditions on the distribution of preference data \cite{bong2022generalized}. Section~\ref{sec:theory} further elaborates on the theoretical underpinnings of DPO and its connections to related literature.

\textbf{Mechanism of the DPO update.} To gain insight into the workings of DPO, it is instructive to examine the gradient of its loss $\mathcal{L}_\text{DPO}$. Taking the derivative with respect to $\theta$ yields:

\begin{multline*}\label{eq:gradient}
    \nabla_\theta \mathcal{L}_\text{DPO}(\pi_\theta;\pi_\text{ref}) = \\ -\beta\mathbb{E}_{(x, y_w, y_l) \sim \mathcal{D}} \bigg[\underbrace{\sigma(\hat{r}_\theta(x, y_l) - \hat{r}_\theta (x, y_w))}_\text{assigns greater emphasis when reward prediction is inaccurate}\bigg[\underbrace{\nabla_\theta\log \pi(y_w \mid x)}_\text{encourage $y_w$} - \underbrace{\nabla_\theta\log\pi(y_l \mid x)}_\text{discourage $y_l$}\bigg]\bigg],
\end{multline*}

where the implicit reward $\hat{r}_\theta(x, y) = \beta \log \frac{\pi_\theta(y \mid x)}{\pi_\text{ref}(y \mid x)}$ is specified by the model $\pi_\theta$ relative to the reference model $\pi_\text{ref}$ (see Section~\ref{sec:theory} for details). Conceptually, the gradient of $\mathcal{L}_\text{DPO}$ increases the model’s probability of generating preferred continuations $y_w$ while reducing the likelihood of the less-preferred $y_l$. The adjustment for each example depends on how much more highly the implicit reward function rates the less-preferred completion $y_l$, with the factor $\beta$ reflecting both the misordering severity and the magnitude of the KL regularization. Empirical studies highlight the importance of this weighting; omitting this coefficient, as in a naive variant of the algorithm, can lead the language model to degenerate behavior (see Appendix Table~\ref{tab:unlikelihood_generations}).

\textbf{DPO overview.} The standard DPO process involves: 1) For each prompt $x$, generating two responses $y_1, y_2$ sampled from $\pi_\text{ref}(\cdot \mid x)$, and using human preference annotations to form an offline preferences dataset $\mathcal{D} = \{x^{(i)}, y_w^{(i)}, y_l^{(i)}\}_{i=1}^N$; 2) Updating the language model $\pi_\theta$ by minimizing the loss $\mathcal{L}_\text{DPO}$ with respect to the fixed reference $\pi_\text{ref}$, dataset $\mathcal{D}$, and a chosen $\beta$.  
In real-world scenarios, existing public preference datasets are often reused in place of collecting new completions and annotations. As these datasets are typically collected using $\pi^\text{SFT}$, it is standard to initialize $\pi_\text{ref}$ with $\pi^\text{SFT}$, where available. If no $\pi^\text{SFT}$ checkpoint is provided, $\pi_\text{ref}$ is instead fit via maximum likelihood estimation to the preferred responses ${(x, y_w)}$ in $\mathcal{D}$, i.e., by setting ${\pi_\text{ref} = \argmax_{\pi}\mathbb{E}_{x, y_w \sim \mathcal{D}}[\log \pi(y_w \mid x)]}$. This step helps alleviate mismatch between the unavailable true reference distribution and the chosen $\pi_\text{ref}$ within DPO training. Additional specifics about hyperparameters and implementation are included in Appendix~\ref{app:implementation}.

\section{Theoretical Analysis of DPO} This section presents deeper insight into the DPO technique, establishes theoretical foundations, and clarifies the benefits of DPO in relation to the drawbacks found in actor-critic approaches frequently applied in RLHF (such as PPO~\cite{schulman2017proximal}).

\label{sec:theory}

\subsection{Your Language Model Is Secretly a Reward Model} DPO achieves policy learning via a single maximum likelihood objective, foregoing both explicit reward modeling and conventional RL algorithms. The loss in Eq. \ref{eq:main_eq} can be viewed as instantiating a Bradley-Terry preference model where the reward is parameterized as $r^*(x, y) = \beta \log\frac{\pi^*_\theta(y \mid x)}{\pi_\text{ref}(y \mid x)}$, with $\pi_{\theta}$ being optimized just as in reward model training (cf. Eq. \ref{eq:reward_model}) under a reparameterization. Here, we formalize the theoretical perspective underlying this parameterization, demonstrate that it does not impose restrictions on the space of attainable reward models, and prove that it can exactly recover the optimal policy. We initiate this discussion by introducing an equivalence relation for reward functions.

\begin{definition}
We define two reward functions $r(x, y)$ and $r'(x, y)$ as equivalent if and only if ${r(x, y)-r'(x, y) = f(x)}$ for some function $f$.
\end{definition}

This relation is easily verified to satisfy the properties of an equivalence relation, thereby dividing the space of reward functions into equivalence classes. We now present two key lemmas:

\begin{lemma}\label{lemma:same_prefrence}
    Within the context of the Plackett-Luce model—specifically, the Bradley-Terry framework—any pair of reward functions that belong to the same equivalence class produce identical preference distributions.
\end{lemma}

\begin{lemma}\label{lemma:same_policy}
    Any two reward functions within the same equivalence class will induce the same optimal policy for the corresponding constrained RL formulation.
\end{lemma}

The arguments are elementary, and their proofs can be found in Appendix \ref{app:lemma1}. The first lemma echoes the well-known non-identifiability problem present in the Plackett-Luce family \cite{plackett1975analysis}. To resolve this, additional identifiability conditions are typically required to ensure meaningful MLE solutions for Eq.~\ref{eq:reward_model} \cite{bong2022generalized}. The second lemma guarantees that the optimal policy is invariant within an equivalence class of reward functions; thus, in our analysis, recovering any representative reward function from the optimal class is sufficient. We formalize this in the theorem below, whose proof is deferred to Appendix~\ref{app:thm1}:

\begin{theorem}\label{thm:main}
    Provided mild regularity conditions hold, every equivalence class of reward functions compatible with the Plackett-Luce (including Bradley-Terry) models may be characterized by the reparameterization ${r(x, y) = \beta \log \frac{\pi(y\mid x)}{\pi_\text{ref}(y\mid x)}}$ for some model $\pi(y\mid x)$ and a fixed reference $\pi_\text{ref}(y \mid x)$.
\end{theorem}

\begin{sproof}
    Let $r(x, y)$ be an arbitrary reward function, with the corresponding optimal policy $\pi_r(y \mid x)$ as defined by Eq.~\ref{eq:op_policy}. Our goal is to show that some member of the equivalence class containing $r$ can be expressed via the stated reparameterization. To this end, we introduce the projection
\begin{equation}
    f(r; \pi_\text{ref}, \beta)(x, y) = r(x, y) - \beta\log\sum_{y}\pi_\text{ref}(y\mid x)\exp\left(\frac{1}{\beta}r(x, y)\right)
\end{equation}
This operator effectively normalizes the reward by subtracting the log-partition function associated with $\pi_r$. Notably, the normalization depends only on $x$, so $f(r; \pi_\text{ref}, \beta)(x, y)$ remains within the same equivalence class as $r(x, y)$. Substituting $r$ with the right-hand side of Eq.~\ref{eq:main_eq}, which holds generally, yields $f(r; \pi_\text{ref}, \beta)(x, y) = \beta \log \frac{\pi_r(y\mid x)}{\pi_\text{ref}(y\mid x)}$. Hence, the projection $f$ constructs an explicit representative of the equivalence class in the desired functional form, establishing that no expressivity is lost under this reparameterization.
\end{sproof}

Theorem~\ref{thm:main} can also be interpreted as identifying, within each equivalence class, the precise reward function that the DPO reparameterization chooses. Specifically, this is the reward function that fulfills the following condition:

\begin{equation}\label{eq:lag_p}
     \sum_{y}\underbrace{\pi_\text{ref}(y\mid x)\exp\left(\frac{1}{\beta}r(x, y)\right)}_{=\pi(y\mid x)\text{, using Thm.~\ref{thm:main} reparam.}} = 1,
\end{equation}

In other words, $\pi(y\mid x)$ forms a proper probability distribution, being both normalized and positive. Furthermore, referencing Eq.~\ref{eq:op_policy}, we observe that Eq.~\ref{eq:lag_p} serves as the partition function corresponding to the optimal policy that arises from the reward function $r(x, y)$. The central innovation behind the DPO algorithm is the enforcement of specific constraints on the otherwise underdetermined Plackett-Luce (and, by extension, Bradley-Terry) class of preference models. This approach retains the expressiveness of the set of possible reward functions, while also ensuring that the optimal policy expressed in Eq.~\ref{eq:op_policy} can be computed in closed form for any given prompt $x$.

\subsection{Instability of Actor-Critic Algorithms} Our framework is equally useful for identifying sources of instability in popular actor-critic methods used in RLHF pipelines, including PPO. To analyze this, we consider the RL fine-tuning stage described in Section \ref{section:prelims}. This analysis allows us to establish links to the control as inference perspective \cite{levine2018reinforcement}, as applied to the constrained RL formulation in \ref{eq:RL}. Here, we parameterize our model as $\pi_{\theta}(y\mid x)$ and aim to minimize the divergence $\mathbb{D}_{\text{KL}}[\pi_{\theta}(y|x) \mid \mid \pi^*(y\mid x)]$, where $\pi^*$ denotes the optimal policy defined by Eq.~\ref{eq:optimum_model}, itself derived from the reward function $r_{\phi}(y, x)$. Carrying out some manipulations leads us to the following objective function:

\begin{equation}\label{eq:AC}
    \max_{\pi_{\theta}}\mathbb{E}_{\pi_{\theta}(y\mid x)}\left[\underbrace{r_{\phi}(x, y) -\beta\log\sum_{y}\pi_\text{ref}(y\mid x)\exp\left(\frac{1}{\beta}r_{\phi}(x, y)\right)}_{f(r_{\phi}, \pi_\text{ref}, \beta)} - \underbrace{\beta\log\frac{\pi_{\theta}(y\mid x)}{\pi_\text{ref}(y\mid x)}}_{\text{KL}}\right]
\end{equation}

This objective is identical to the one targeted in earlier research 
\citep{ziegler2020finetuning, stiennon2022learning, bai2022training, ouyang2022training}, where the reward class $r_{\phi}$ employs the DPO-equivalent reward. Within this framework, the normalization component in $f(r_{\phi}, \pi_\text{ref}, \beta)$ can be viewed as the soft value function corresponding to the reference policy $\pi_\text{ref}$. Although this normalization does not alter the set of optimal solutions, its absence can cause the policy gradient estimator to exhibit significant variance, thereby destabilizing the learning process. One way to address this normalization is to employ a trainable value function, though this approach itself is often challenging to optimize. Another method used in previous studies is to normalize rewards using a human completion as a baseline, effectively a single-sample Monte-Carlo estimate of the normalizing constant. The DPO reparameterization, however, produces a reward function formulation that eliminates the need for any baseline.

\section{Experiments}
This section presents empirical analyses of DPO’s performance in training policies solely from preference data. Initially, we investigate a controlled text-generation environment to assess how effectively DPO negotiates the balance between reward maximization and minimizing KL-divergence with the reference policy, in comparison to widely-used preference-based algorithms such as PPO. We then extend our examination to encompass larger models and more challenging RLHF applications, including tasks like summarization and dialogue. Our results show that with minimal hyperparameter adjustment, DPO generally matches or surpasses robust baselines like RLHF with PPO, and even outperforms strategies that select the best of $N$ trajectories as ranked by a learned reward model. A thorough account of our experimental methodology precedes these findings, with further experimental specifics provided in Appendix~\ref{app:exp_details}.

\textbf{Tasks.} Our empirical analysis encompasses three distinct open-ended text generation scenarios. Across all studies, learning algorithms infer a policy using a collection of preference-labeled data $\mathcal{D}=\bigl\{x^{(i)}, y_w^{(i)}, y_l^{(i)}\bigr\}_{i=1}^N$. In the \textbf{controlled sentiment generation} setting, $x$ represents a partial movie review excerpted from the IMDb dataset \cite{maas-EtAl:2011:ACL-HLT2011}, and the objective is for the policy to produce a continuation $y$ reflecting positive sentiment. To facilitate systematic evaluation, we construct preference pairs automatically using a pre-trained sentiment classifier, selecting those for which $p(\text{positive}\mid x,y_w)>p(\text{positive}\mid x,y_l)$. For SFT, we adapt GPT-2-large through supervised fine-tuning on the training subset of IMDB reviews, running until convergence (see App~\ref{app:sentiment_details} for additional specifics). In the \textbf{summarization} context, $x$ corresponds to a Reddit forum post, and the policy is required to produce a summary $y$ highlighting its main points. Consistent with established benchmarks, we leverage the Reddit TL;DR summarization dataset \citep{volske-etal-2017-tl} together with human preference judgments collected by \citeauthor{stiennon2022learning}. An SFT model fine-tuned on reference summaries of forum posts\footnote{\url{https://huggingface.co/CarperAI/openai_summarize_tldr_sft}} is trained using the TRLX \citep{leandro_von_werra_2023_7790115} RLHF framework. Human preferences were originally obtained by \citeauthor{stiennon2022learning} using generations from a similar but independently fine-tuned SFT model. For the \textbf{single-turn dialogue} task, $x$ is a user-posed query, potentially ranging from scientific questions to requests for personal guidance. The policy is tasked with producing a helpful and engaging reply $y$; experiments utilize the Anthropic Helpful and Harmless dialogue dataset \citep{bai2022training}, containing 170,000 human–assistant conversational turns. Each dialogue culminates with two responses produced by a large, though unspecified, language model, each labeled to indicate human preference. Since there is no available pre-trained SFT model in this case, we construct an SFT baseline by fine-tuning a standard language model exclusively on the preferred responses.

\textbf{Evaluation.} We employ two main strategies to evaluate our approaches. To precisely measure how well each algorithm performs on the objective of constrained reward maximization, we assess their performance in the controlled sentiment generation scenario by plotting the trade-off curve, or frontier, between achieved reward and KL-divergence relative to the reference policy; this is possible because we have direct access to the true reward signal, operationalized as a sentiment classifier. In practical applications, however, the true reward function is not available. Accordingly, for summarization and single-turn dialogue, we measure each algorithm’s \textit{win rate} against a baseline policy, utilizing GPT-4 to approximate human preferences regarding summary quality and dialogue helpfulness, respectively. For the summarization task, the baseline consists of reference summaries from the test set; for dialogue, it is the human-preferred reply provided in the test set. While prior research indicates that language models may surpass traditional automatic metrics as evaluators \citep{Chen2023ExploringTU}, we validate our reliance on GPT-4 by conducting a human evaluation as described in Sec.~\ref{sec:human-judgments}. Our findings demonstrate a strong correlation between GPT-4 assessments and human judgments, with levels of human–GPT-4 agreement comparable to or surpassing inter-annotator consistency among humans. 

\textbf{Methods.} Beyond DPO, we benchmark a variety of existing techniques for aligning language model outputs with human preference. For the summarization task, we test zero-shot prompting of \textbf{GPT-J} \citep{gpt-j}, while for dialogue, we employ 2-shot prompting with \textbf{Pythia-2.8B} \citep{biderman2023pythia}. Additionally, our comparisons include the standard \textbf{SFT} model as well as \textbf{Preferred-FT}, where the latter is produced by fine-tuning via supervised learning using the chosen responses $y_w$—sourced from SFT in the sentiment and summarization settings and from a standard LM in dialogue. We also investigate the pseudo-supervised \textbf{Unlikelihood} approach~\citep{welleck2019neural}, which updates the model to increase the likelihood of $y_w$ and decrease it for $y_l$, employing a tunable coefficient $\alpha \in [0,1]$ to weigh the ‘unlikelihood’ component. Furthermore, we include \textbf{PPO} \citep{schulman2017proximal}, which trains on a reward function learned from preference data, as well as \textbf{PPO-GT}, an oracle variant that leverages access to the ground-truth reward in the sentiment-controlled scenario. For the sentiment experiments, we evaluate two PPO-GT implementations: an off-the-shelf version \cite{leandro_von_werra_2023_7790115} and a modified version that applies reward normalization and refined hyperparameter tuning for enhanced outcomes (similar adjustments are applied to PPO with learned rewards). Lastly, we incorporate the \textbf{Best of $N$} method, where $N$ completions are generated from either the SFT model (or Preferred-FT in dialogue) and the response with the highest learned reward score is chosen. While this strategy can yield strong results by isolating reward model performance from PPO optimization, it remains computationally expensive even for moderate values of $N$ due to the need to sample multiple completions per input during evaluation.

\subsection{How well can DPO optimize the RLHF objective?}

In standard RLHF approaches, the objective of maximizing reward is moderated by a KL-divergence constraint, ensuring that the updated policy does not stray too far from the reference policy. As a result, algorithmic comparisons should jointly consider both attained reward and the KL divergence from the reference policy; maximizing reward at the expense of a much larger KL may not constitute an improvement. Figure~\ref{fig:frontier-tldr-main} presents the reward versus KL tradeoff for different algorithms within the sentiment classification task. For each algorithm, we conduct several training runs, systematically varying the hyperparameter associated with policy conservativeness—specifically, setting the target KL to $\{3,6,9,12\}$ for PPO, $\beta$ to $\{0.05,0.1,1,5\}$ for DPO, $\alpha$ to $\{0.05,0.1,0.5,1\}$ for the unlikelihood method, and employing different random seeds for preferred-FT. In total, 22 separate training runs are included. At intervals of every 100 training steps up to convergence, policies are evaluated on a fixed set of test prompts by computing the average reward under the authentic reward function and determining the average sequence-level KL\footnote{Here, the sequence-level KL refers to the aggregate of per-timestep KL divergences.} relative to the reference policy, $\text{KL}\left(\pi\mid \mid \pi_\text{ref}\right)$. The results indicate that DPO generates a reward-KL frontier that is markedly more favorable than the others, attaining the highest reward without incurring high KL divergence. This finding is particularly significant for several reasons: First, although both DPO and PPO aim to optimize the identical objective, DPO exhibits substantially greater efficiency, with its reward/KL Pareto frontier strictly surpassing that of PPO. Second, DPO maintains a superior reward-KL frontier compared to PPO, \emph{even in scenarios where PPO is provided access to ground truth rewards} (PPO-GT).

\subsection{Can DPO scale to real preference datasets?}
\label{sec:dpo-real-datasets}
We proceed by investigating how DPO performs when fine-tuned on real-world preference data, specifically in the domains of summarization and single-turn dialogue. Notably, standard automated metrics like ROUGE often exhibit weak alignment with actual human judgments in summarization tasks~\citep{stiennon2022learning}. Previous studies have demonstrated that models fine-tuned via PPO, based on human preference signals, tend to generate higher-quality summaries. To assess various approaches, we generate completions on the test partition of the TL;DR summarization dataset and measure the mean win rate when compared against reference outputs from the same set. For each approach, outputs are sampled over a temperature range from 0.0 to 1.0, with win rates depicted in Figure~\ref{fig:frontier-tldr-main} (right). DPO, PPO, and Preferred-FT are all applied to further train the identical GPT-J SFT model\footnote{\url{https://huggingface.co/CarperAI/openai_summarize_tldr_sft}}. In our evaluation, DPO attains a win rate close to 61\% at a temperature of 0.0, outperforming PPO, which peaks around 57\% at the same temperature. In addition, DPO secures a superior maximum win rate compared to the best result from the $N$ baseline. It is important to mention that we did not perform extensive tuning of DPO’s $\beta$ hyperparameter, and thus the results presented may not fully reflect the method’s capabilities. Furthermore, DPO demonstrates a greater resilience to temperature variations, in contrast with PPO, whose performance can deteriorate to that of the original GPT-J at elevated temperatures. Preferred-FT fails to yield notable gains over the initial SFT model. Direct human assessments comparing DPO and PPO, reported in Section~\ref{sec:human-judgments}, reveal that samples from DPO at temperature 0.25 were preferred 58\% of the time over PPO samples at the same temperature.

For single-turn dialogue, we assess various approaches using a subset of the test split from the Anthropic HH dataset \citep{bai2022training}, focusing on cases involving a single round of human-assistant interaction. GPT-4 evaluations utilize the preferred completions from the test set as ground truth to determine win rates across methods. In the absence of an established SFT model for this benchmark, our workflow begins with a pre-trained Pythia-2.8B, applies Preferred-FT to fine-tune a reference model on preferred completions—ensuring output distributional alignment—and subsequently trains with DPO. For comparison, we report results from the strongest of 128 Preferred-FT completions (having observed that the Best of $N$ baseline saturates at $N=128$ for this scenario; see Appendix Figure~\ref{fig:best-of-n}) as well as a 2-shot prompt-based version of the Pythia-2.8B model. DPO achieves equal or superior performance compared to these baselines when using optimal temperatures for each method. Additionally, we assess an RLHF model employing PPO trained on the Anthropic HH dataset \footnote{\url{https://huggingface.co/reciprocate/ppo_hh_pythia-6B}} from a reputable repository \footnote{\url{https://github.com/CarperAI/trlx/tree/main/examples/hh}}, but are unable to identify a prompt or temperature that outperforms the base Pythia-2.8B. Given our TL;DR findings and that both approaches maximize the same reward signal, we take the Best of 128 baseline as a reasonable approximation of PPO-level performance. In sum, DPO stands out as the only computationally efficient method to surpass the preferred completions from Anthropic HH, yielding results comparable to or better than the computationally intensive Best of 128. As indicated by Figure~\ref{fig:dialogue-main}, DPO also reaches its optimal performance in relatively few training steps.

\subsection{Generalization to a new input distribution}

To further investigate the robustness of PPO and DPO under distributional changes, we assess the policies trained via PPO and DPO from our Reddit TL;DR summarization setup on a distinct input: the test split of news articles from the CNN/DailyMail dataset \citep{nallapati-etal-2016-abstractive}. For evaluation, we apply the optimal sampling temperatures determined from the TL;DR task (0 and 0.25). Table~\ref{tab:ood} summarizes the findings. We calculate the GPT-4 win rate relative to the reference summaries in each dataset, utilizing the same GPT-4 (C) prompt previously employed for Reddit TL;DR, but substituting the phrase ``forum post'' with ``news article''. Across this new data distribution, DPO maintains a strong lead over the PPO policy. These observations suggest that DPO-trained policies retain strong generalization capabilities, on par with those learned through PPO—even though DPO is not provided the additional unlabeled Reddit TL;DR prompts leveraged by PPO.

\subsection{Validating GPT-4 judgments with human judgments}
\label{sec:human-judgments}
We implement a human evaluation to assess the trustworthiness of GPT-4’s decisions, leveraging results from the TL;DR summarization experiments and two distinct GPT-4 prompts. The \textbf{GPT-4 (S)} (simple) prompt solicits which summary more effectively conveys the essential content of the post. The \textbf{GPT-4 (C)} (concise) prompt additionally requests a judgment of conciseness; this alternative is examined since, with the \textbf{GPT-4 (S)} prompt, GPT-4 tends to select lengthier, more repetitive summaries compared to human preferences. Full prompt wording is detailed in Appendix~\ref{app:prompts}. For a broad assessment of sample quality, we conduct three types of comparisons: using the best-performing model (DPO, temp. 0.25), the lowest-performing one (PPO, temp. 1.0), and an intermediate model (SFT, temp. 0.25). All three are evaluated against PPO with greedy sampling at its optimal temperature. Analysis reveals that for both prompts, GPT-4’s selections align with human choices about as often as human raters agree with each other, indicating that GPT-4 serves as a reliable stand-in for human evaluation (given a limited pool of human annotators, repeated human ratings are obtained only for the DPO and PPO-1 scenarios). Notably, win rates using the \textbf{GPT-4 (C)} prompt more closely track those observed with human judgments; consequently, we employ this prompt as the principal metric in Section~\ref{sec:dpo-real-datasets}. Further procedural details—such as the user interface for raters and the roster of participating human annotators—are provided in Appendix~\ref{app:human-study}.

\section{Discussion}
Preference-based learning represents a robust and scalable approach for developing proficient and well-aligned language models. In this work, we introduced DPO, a straightforward method for leveraging preferences to train language models that circumvents reinforcement learning. Instead of reinterpreting preference learning as a standard RL problem to accommodate existing RL solutions, DPO establishes a correspondence between policies of language models and reward functions. This allows for direct optimization of human preferences through a basic cross-entropy loss, eliminating the need for reinforcement learning and without sacrificing expressiveness. Notably, DPO achieves competitive or superior outcomes compared to prevalent RLHF algorithms such as those utilizing PPO, and requires minimal hyperparameter adjustment, thus significantly lowering the complexity of preference-based language model training.

\textbf{Limitations \& Future Work.} Our findings highlight several directions for continued research. One important area is to understand how DPO-trained policies generalize outside the observed data compared to models trained using explicit reward functions; initial comparisons indicate DPO performs on par with PPO-based approaches, though broader investigations are necessary. Further, it remains an open question whether self-labeling using DPO can facilitate effective use of prompts without supervision. Another line of inquiry involves examining how over-optimization of the implicit reward in DPO manifests, and whether the modest performance drop noted in Figure~\ref{fig:dialogue-main}-right can be attributed to this phenomenon. Moreover, while the present work covers models containing up to 6B parameters, scaling DPO for models that are orders of magnitude larger is a promising path for future study. Regarding evaluation protocols, our observations suggest that GPT-4-derived win rates are sensitive to prompt formulations; future exploration into eliciting more reliable judgments from automated systems is warranted. Ultimately, DPO has potential utility that extends beyond the domain of human-preference-based language model training, potentially benefiting the training of generative models in diverse modalities.